\section{排列 组合问题(先不要看这部分)}

\textbf{为了便于了解图的存储结构,这两个题目用的是邻接表来解决的，所以这部分可以放到稍后再看。}

\subsection{P1706 全排列问题}
\begin{lstlisting}[language=c++]
#include<bits/stdc++.h>
using namespace std;
constexpr int N = 11;
vector<int> e[N];
int n;
int path[N];
bool used[N];

void dfs(const int &u, const int &depth) {
	if (depth == n) {
		for (int i = 0; i < n; i++) {
			cout << setw(5) << path[i];
		}
		cout << endl;
		return;
	}
	for (const auto &v: e[u]) {
		if (used[v]) {
			continue;
		}
		used[v] = true;
		path[depth] = v;
		dfs(v, depth + 1);
		used[v] = false;
	}
}

int main() {
	cin >> n;
	for (int i = 0; i <= n; i++) {
		for (int v = 1; v <= n; v++) {
			e[i].emplace_back(v);
		}
	}
	dfs(0, 0);
	return 0;
}
\end{lstlisting}

\subsection{P1157 组合的输出}
\begin{lstlisting}[language=c++]
#include <bits/stdc++.h>
using namespace std;

constexpr int N = 21;
vector<int> e[N];
int n, r;
int path[N];

void dfs(const int &u, const int &depth) {
	if (depth == r) {
		for (int i = 0; i < r; i++) {
			cout << setw(3) << path[i];
		}
		cout << '\n';
		return;
	}
	
	for (const auto &v: e[u]) {
		path[depth] = v;
		//cout << "u:" << u << " v:" << v << '\n'; //递归过程
		dfs(v, depth + 1);
	}
}

int main() {
	cin >> n >> r;
	for (int i = 0; i <= n; i++) {
		for (int v = i + 1; v <= n; v++) {
			e[i].push_back(v);
		}
	}
	
	dfs(0, 0);
	return 0;
}

\end{lstlisting}